\newcommand{\sfL}{\mathsf{L}}
\newcommand{\sfS}{\mathsf{S}}
\newcommand{\sfT}{\mathsf{T}}
\newcommand{\clG}{\mathcal{G}}
\newcommand{\clV}{\mathcal{V}}
\newcommand{\clZ}{\mathcal{Z}}
\newcommand{\EE}{\mathbb{E}}
\newcommand{\RR}{\mathbb{R}}
\newcommand{\Frob}{\mathrm{F}}
\newcommand{\dif}{\mathop{}\!\mathrm{d}}
\newcommand{\ttR}{\mathtt{R}}
\newcommand{\ttB}{\mathtt{B}}
\newcommand{\triplenorm}[1]{\lVert\!\lVert\!\lVert #1\rVert\!\rVert\!\rVert}

\section{Main Results}

\subsection{The Approximate SPDE}

\begin{proposition}
    Let $\sigma \in \text{Lip}(\R, (0,\infty))$, and $\ell\in(2,4]$.There is a constant $C = C(\sigma,\ell) < \infty$ such the following holds. Let $v_t^{\rho}$ solve the following equation with $v_0^{\rho} \in \mathscr{X}_0^{\ell}$
    \[
    v_t^{\rho}(x) = \mathcal{G}_{t-s}v_s^{\rho}(x) + \gamma_{\rho}\int_s^t\mathcal{G}_{t+\rho-r}[\sigma_{\rho}(v_r^{\rho})dW_r](x)
    \]
    where
    \[
    \sigma_{\rho}(x) := \sigma(x/\gamma_{\rho}),\quad x\in \R.
    \]
    For $T_0,T_1$ satisfying
    \[
    (TBA)
    \]
    we define for $t\geq T_0$,
    \[
    \widehat{B}_t^{\rho}(x) = \gamma_{\rho}\int_{T_0}^t \mathcal{G}_{T_1+\rho-r}[d\widetilde{W}_r^{\rho}](x)
    \]
    where $d\widetilde{W}^{\rho}$ satisfying
    \[
    \left\lVert \int_{T_0}^t\mathcal{G}_{\eta_r}[\sigma_{\rho}(v_r^{\rho})dW_r - (\mathcal{S}_{\zeta_r}\sigma_{\rho}^2(v_r^{\rho}))^{1/2}d\widetilde{W}_r]^{\rho}\right\rVert_2^2 \leq \dfrac{1}{2\pi^3}\int_{T_0}^t \zeta_r^2\eta_r^{-2}\lVert|\sigma_{\rho}(v_r^{\rho})|\rVert_2^2dr
    \]
    for $\zeta_r = [L(1/\rho)^{-1}(T_1-r)]^{1/2}$ and arbitrary postive path $(\eta_r)_{r\geq T_0}$, where
    \[
    \lVert|\sigma_{\rho}(v_r^{\rho})|\rVert_l^l :=\sup_{x\in\R^2}\mathbb{E}|\sigma_{\rho}(v_r^{\rho})(x)|^l.
    \]
    Then we have
    \[\begin{aligned}
\mathbb{E}
&\left|
V_t^{\rho,T_1}(x)
-
V_{T_0}^{\rho,T_1}(x)
-
\int_{T_0}^{t}
J_{\sigma_{\rho}}\left(
    S_{\rho}(T_1-r),
    V_r^{\rho,T_1}(x)
\right)
\,d\widehat{B}_r(x)
\right|^2 \\
\leq
&C\left\langle \lVert v_0\rVert_{\ell}\right\rangle^2
\frac{L\!\left(L(1/\rho)\right)}{L(1/\rho)}
\end{aligned}
\]
where
\[
V_s^{\rho,t}(x):= \mathcal{G}_{t-s}v_s^{\rho}(x).
\]
\end{proposition}
\begin{proof}
    We may notice that, for all
\(t \in [T_0,T_1]\),
\begin{equation*}
V_t^{\rho,T_1}(x)
=
V_{T_0}^{\rho,T_1}(x)
+
\gamma_\rho
\int_{T_0}^{t}
\mathcal{G}_{T_1+\rho-r}
\bigl[\sigma_{\rho}(v_r^{\rho})\,dW_r\bigr](x).
\end{equation*}
So that we know
\[
\begin{aligned}
    &\left(\mathbb{E}
\left|
V_t^{\rho,T_1}(x)
-
V_{T_0}^{\rho,T_1}(x)
-
\int_{T_0}^{t}
J_{\sigma_{\rho}}\left(
    S_{\rho}(T_1-r),
    V_r^{\rho,T_1}(x)
\right)
\,d\widehat{B}_r(x)
\right|^2\right)^{1/2} \\
&\leq 
\left(\mathbb{E}
\left|
\gamma_\rho
\int_{T_0}^{t}
\mathcal{G}_{T_1+\rho-r}
\left[
 \sigma_{\rho}(v_r^{\rho})\,dW_r
 -
 \left(\mathcal{S}_{\zeta_r}
       [\sigma_{\rho}^2\circ v_r^{\rho}]\right)^{1/2}
 d\widetilde W_r^{\rho}
\right](x)
\right|^2\right)^{1/2} + \\
&\qquad\left(\mathbb{E}\left|
\gamma_{\rho}\int_{T_0}^t \mathcal{G}_{T_1+\rho-r}\left[\left(\mathcal{S}_{\zeta_r}[\sigma_{\rho}^2\circ v_r^{\rho}]^{1/2}\right)d\widetilde{W}_r^{\rho}\right](x)\right.\right.\\ &\qquad- \left.\left.\int_{T_0}^t J_{\sigma_{\rho}}(S_{\rho}(T_1-r),V_r^{\rho,T_1}(x))d\widehat{B}_r(x) 
\right|^2\right)^{1/2}
\end{aligned}
\]
We choose
\(\eta_r := T_1+\rho-r\) to find a
space-time white noise \(d\widetilde W^{\rho}\) such that, for any
\(x\in\mathbb{R}^2\),
\begin{align*}
&\mathbb{E}
\left|
\gamma_\rho
\int_{T_0}^{t}
\mathcal{G}_{T_1+\rho-r}
\left[
 \sigma_{\rho}(v_r)\,dW_r
 -
 \left(\mathcal{S}_{\zeta_r}
       [\sigma_{\rho}^2\circ v_r^{\rho}]\right)^{1/2}
 d\widetilde W_r^{\rho}
\right](x)
\right|^2
\nonumber\\
&\qquad\le
\dfrac{2}{\pi^2L(1/\rho)^2}\int_{T_0}^t\dfrac{T_1-r}{(T_1+\rho-r)^2}\lVert|\sigma_{\rho}(v_r^{\rho})|\rVert_2^2 dr \\
&\qquad\le
\dfrac{2\lVert \sigma\rVert_{\infty}^2}{\pi^2L(1/\rho)^2}\int_{T_0}^{T_1} \dfrac{T_1-r}{(T_1+\rho-r)^2}dr \\
&\qquad\le
\dfrac{2\lVert \sigma\rVert_{\infty}^2}{\pi^2L(1/\rho)^2}\left(L_{\rho}(T_1-T_0) + 1 - \rho/(T_1-T_0 +\rho)\right)
\leq \dfrac{C(\sigma)S_{\rho}(T_1-T_0)}{L(1/\rho)}
\end{align*}

Then, we compute
\[
\begin{aligned}
&\mathbb{E}\left|
\gamma_{\rho}\int_{T_0}^t \mathcal{G}_{T_1+\rho-r}\left[\left(\mathcal{S}_{\zeta_r}[\sigma_{\rho}^2\circ v_r^{\rho}]^{1/2}\right)d\widetilde{W}_r^{\rho}\right](x)\right.\\ &\qquad- \left.\int_{T_0}^t J_{\sigma_{\rho}}(S_{\rho}(T_1-r),V_r^{\rho,T_1}(x))d\widehat{B}_r(x) 
\right|^2 \\
= &\dfrac{4\pi}{L(1/\rho)}\mathbb{E}\left\lvert\int_{T_0}^t \int G_{T_1+\rho-r}(x-y)\left(\mathcal{S}_{\zeta_r}[\sigma_{\rho}^2\circ v_r^{\rho}]^{1/2}\right)d\widetilde{W}_r^{\rho}(y) \right. \\
&\qquad -\int_{T_0}^t\int J_{\sigma_{\rho}}(S_{\rho}(T_1-r),V_r^{\rho,T_1}(x))G_{T_1+\rho-r}(x-y)d\widetilde{W}_r^{\rho}(y) 
\left.\right\rvert^2 \\
\leq& \dfrac{4\pi}{L(1/\rho)}\int_{T_0}^t\int G_{T_1+\rho-r}(x-y)^2F(x,y)dydr
\end{aligned}
\]
where
\[
F(x,y):= \mathbb{E}\left|\left(\mathcal{S}_{\zeta_r}[\sigma_{\rho}^2\circ v_r^{\rho}]^{1/2}\right)(y) - J_{\sigma_{\rho}}(S_{\rho}(T_1-r),V_r^{\rho,T_1}(x))\right|^2
\]
\end{proof}

\subsection{First stage approximation}

\begin{proposition}
    Fix \(M,\lambda\in(0,\infty)\), \(\ell\in(2,4]\),
\(\overline{\beta}>\eta(\ell)\), where
\[
\eta(\ell)
:=
\left[
2^{-2/\ell}(\ell-1)
\left(\ell-{2}\right)
\right]^{1/2}.
\]
There is a constant
\[
C=C(M,\lambda,\ell,\overline{\beta})
\in(0,\infty)
\]
such that the following holds.
Suppose that \(\rho\in(0,C^{-1})\).
Let \(\sigma\) satisfy
\(\operatorname{Lip}(\sigma)\leq\lambda\), $|\sigma|_{\infty}^2 \leq M$, and let
\((v_t)_{t\geq 0}\) solve
\[
    v_t^{\rho}(x) = \mathcal{G}_{t-s}v_s^{\rho}(x) + \gamma_{\rho}\int_s^t\mathcal{G}_{t+\rho-r}[\sigma_{\rho}(v_r^{\rho})dW_r](x)
\]
with initial condition
\(v_0\in\mathscr{X}_0^\ell\).
Let
\(\xi\in[0,(TBA)]\),
\begin{equation*}
t\in
\left[
 \left(
   [1+\mathrm{L}(1/\rho)]\kappa_{\ell,\rho}^{-1}
   +\kappa_{\ell,\rho}
 \right)\xi^2,\, \alpha\rho
\right],\quad \kappa_{\ell,\rho} = L(1/\rho)^{-1/(1-2/\ell)}
\end{equation*}
\textit{ for some $\alpha:=\alpha(\sigma)$ constant, 
\(x,z\in\mathbb{R}^2\), and \(\iota_1,\iota_2\in[0,1]\).
Then we have}
\[
\begin{aligned}
\mathbb{E}&
\left|
 \left(
   \mathcal{S}_{\xi,z}
   [\sigma_{\rho}^2\circ v^{\rho}_t](x)
 \right)^{1/2}
 -
 J_{\sigma_{\rho},\rho}
 \left(
   S_{\rho}\!\left(
     \mathrm{L}(1/\rho)^{l_1}\xi^2
   \right),
   \mathcal{G}_{\mathrm{L}(1/\rho)^{l_2}\xi^2}
   v_t^{\rho}(x)
 \right)
\right|^{l_2} \\
&\quad\leq
C
\left\langle
 \lVert v_0\rVert_{\ell}
\right\rangle^2
\frac{
 \mathrm{L}\!\left(\mathrm{L}(1/\rho)\right)
}{
 \mathrm{L}(1/\rho)
}.
\end{aligned}
\]
\end{proposition}
\begin{proof}
The proof combines the earlier results of this section with specific choices of
parameters. We proceed in several steps.

\medskip
\noindent\textit{Step 1: Choosing the parameters.}
Let
\begin{equation*}
\ttR=\ttB_{\xi,k_{\xi,z}(x)+z},\quad
\tau_0=t-(\kappa_{\ell,\rho}^{-1}+\kappa_{\ell,\rho})\xi^2,
\quad
\tau_1=\tau_0+\kappa_{\ell,\rho}^{-1}\xi^2,
\quad
\tau_2=\tau_1+\kappa_{\ell,\rho}\xi^2=t.
\end{equation*}
Let $X\in\RR^d$ denote the center of $\ttR$. Note that
\begin{equation*}
\tau_2\geq\tau_1\geq\tau_0
\geq\sfL(1/\rho)\kappa_{\ell,\rho}^{-1}\xi^2\geq0
\end{equation*}
by the assumption on $t$ in (5.50).

\medskip
\noindent\textit{Step 2: Flattening the solution.}
By the matrix-valued reverse triangle inequality and the flattening estimate,
there is an absolute constant
$C<\infty$ such that, as long as $\rho$ is small enough
\begin{align*}
E_1
&:= \EE\left\lvert
\left(\fint_{\ttR}\sigma_{\rho}^2(v^{\rho}_t(y))\,\dif y\right)^{1/2}
-\left(\fint_{\ttR} \sigma_{\rho}^2\!\left(
\clV_{\tau_1,\tau_2}\clZ_X\clG_{\tau_1-\tau_0}v^{\rho}_{\tau_0}(y)
\right)\,\dif y\right)^{1/2}
\right\rvert^{2}
\\
&\leq \fint_{\ttR}\EE\left\lvert
\sigma_{\rho}\!\left(\clV_{\tau_0,\tau_2}v^{\rho}_{\tau_0}(y)\right)
-\sigma_{\rho}\!\left(
\clV_{\tau_1,\tau_2}\clZ_X\clG_{\tau_1-\tau_0}v^{\rho}_{\tau_0}(y)
\right)
\right\rvert^{2}\,\dif y
\\&\leq \dfrac{\lambda^2L(1/\rho)}{4\pi}\fint_{\ttR}\EE\left\lvert
\clV_{\tau_0,\tau_2}v^{\rho}_{\tau_0}(y)
-\clV_{\tau_1,\tau_2}\clZ_X\clG_{\tau_1-\tau_0}v^{\rho}_{\tau_0}(y)
\right\rvert^{2}\,\dif y
\\
&\leq C K_{\lambda,\rho,\tau_2-\tau_1}
\Xi_{2;\tau_0,\tau_1}[v_{\tau_0}]^2
\left[
\sfL_\rho\!\left(\frac{\tau_1-\tau_0}{\tau_2-\tau_1+\rho}\right)
+\frac{(1+\lambda^2\delta)(\tau_2-\tau_1)+\xi^2}{\tau_1-\tau_0}
+\delta
\right],
\end{align*}
where
\[
\delta=K_{\lambda,\rho,\tau_2-\tau_1}\sfL(1/\rho)^{-1}.
\]
Substituting the parameters from \eqref{eq:choosexietc}, we get
\[
\delta=K_{\lambda,\rho,\kappa_{\ell,\rho}\xi^2}\sfL(1/\rho)^{-1}
\leq K_{\lambda,\rho,\xi^2}\sfL(1/\rho)^{-1},
\]
and hence
\begin{align}
E_1
&\leq C K_{\lambda,\rho,\xi^2}
\Xi_{2;\tau_0,\tau_1}[v_{\tau_0}]^2
\left[
\sfL_\rho(\kappa_{\ell,\rho}^{-2})
+\kappa_{\ell,\rho}\bigl[(1+\lambda^2\delta)\kappa_{\ell,\rho}+1\bigr]
+\frac{K_{\lambda,\rho,\xi^2}}{\sfL(1/\rho)}
\right].
\tag{5.54}\label{eq:E1bd}
\end{align}
From \eqref{eq:E1bd} and the definitions of $\kappa_{\ell,\rho}$ and
$K_{\lambda,\rho,\xi^2}$, it is clear that if $\rho$ is sufficiently small
(depending on $\overline q$ and $\lambda$), then
\begin{equation}
E_1\leq C\Xi_{2;\tau_0,\tau_1}[v_{\tau_0}]^2
\frac{\sfL(\sfL(1/\rho))}{\sfL(1/\rho)},
\tag{5.55}\label{eq:E1OK}
\end{equation}
for a constant $C$ depending only on $\overline q$ and $\lambda$.

\medskip
\noindent\textit{Step 3: Concentration.}
Define
\[
E_2\coloneqq\EE\left\lvert
\left(\fint_{\ttR}\sigma^2\!\left(
\clV_{\tau_1,\tau_2}\clZ_X\clG_{\tau_1-\tau_0}v_{\tau_0}(y)
\right)\,\dif y\right)^{1/2}
-J_{\sigma,\rho}\!\left(
\sfS_\rho(\tau_2-\tau_1),
\clG_{\tau_1-\tau_0}v_{\tau_0}(x)
\right)
\right\rvert_{\Frob}^{2}.
\]
We consider two cases.

If $\tau_2-\tau_1=\kappa_{\ell,\rho}\xi^2\geq\rho$, then the law-of-large-numbers estimate,
applied with $s=\tau_1$, $t=\tau_2$, and
$a=\clG_{\tau_1-\tau_0}v_{\tau_0}(X)$, gives
\begin{align*}
E_2
&\leq C\Xi_{\ell;\tau_1,\tau_2}
[\clG_{\tau_1-\tau_0}v_{\tau_0}(X)]^2
\left(
\frac{\log[\xi(\tau_2-\tau_1)^{-1/2}]}
{\xi(\tau_2-\tau_1)^{-1/2}}
\right)^{2-4/\ell}
\\
&=C\Xi_{\ell;\kappa_{\ell,\rho}\xi^2}
[\clG_{\tau_1-\tau_0}v_{\tau_0}(X)]^2
\left(
\kappa_{\ell,\rho}^{1/2}
\log[\kappa_{\ell,\rho}^{-1/2}]
\right)^{2-4/\ell}.
\end{align*}
Using the definition of $\kappa_{\ell,\rho}$, there is a constant
$C=C(\ell)<\infty$ such that, if $\rho$ is sufficiently small,
\begin{align}
E_2
&\leq C\Xi_{\ell;\kappa_{\ell,\rho}\xi^2}
[\clG_{\tau_1-\tau_0}v_{\tau_0}(x)]^2
\frac{\sfL(\sfL(1/\rho))^{2-4/\ell}}{\sfL(1/\rho)}
\nonumber\\
&\leq C\Xi_{\ell;\kappa_{\ell,\rho}\xi^2}
[\clG_{\tau_1-\tau_0}v_{\tau_0}(x)]^2
\frac{\sfL(\sfL(1/\rho))}{\sfL(1/\rho)}.
\tag{5.56}\label{eq:E2firstbound}
\end{align}
In the last inequality we used $\ell\leq4$.

If $\tau_2-\tau_1=\kappa_{\ell,\rho}\xi^2<\rho$, then the small-time estimate,
applied with the same parameters, gives
\begin{align*}
E_2
&\leq4\Xi_{2;\tau_1,\tau_2}
[\clG_{\tau_1-\tau_0}v_{\tau_0}(X)]^2
\lambda^2\sfS_\rho(\rho)
\\
&=4\Xi_{2;\kappa_{\ell,\rho}\xi^2}
[\clG_{\tau_1-\tau_0}v_{\tau_0}(X)]^2
\lambda^2\frac{\sfL(2)}{\sfL(1/\rho)}.
\end{align*}
Thus, in either case,
\begin{equation}
E_2\leq C\Xi_{\ell;\tau_1,\tau_2}
[\clG_{\tau_1-\tau_0}v_{\tau_0}(X)]^2
\frac{\sfL(\sfL(1/\rho))}{\sfL(1/\rho)},
\tag{5.57}\label{eq:E2finalbd}
\end{equation}
for a constant $C<\infty$ depending only on $\ell$ and $\lambda$.

\medskip
\noindent\textit{Step 4: Adjusting the arguments of $J_{\sigma,\rho}$.}
Let $\iota_1,\iota_2\in[0,1]$. By the regularity estimate for
$J_{\sigma,\rho}$ and the triangle inequality
on $L^2(\Omega,\mathbb P)$,
\begin{align}
&\left(\EE\left\lvert
J_{\sigma,\rho}\!\left(
\sfS_\rho(\kappa_{\ell,\rho}\xi^2),
\clG_{\kappa_{\ell,\rho}^{-1}\xi^2}v_{\tau_0}(X)
\right)
-J_{\sigma,\rho}\!\left(
\sfS_\rho(\sfL(1/\rho)^{\iota_1}\xi^2),
\clG_{\sfL(1/\rho)^{\iota_2}\xi^2}v_t(X)
\right)
\right\rvert_{\Frob}^2\right)^{1/2}
\nonumber\\
&\quad\leq\lambda K_{\lambda,\rho,\xi^2}^{1/2}
\Biggl[
\Xi_{2;\sfL(1/\rho)\xi^2}
[\clG_{\kappa_{\ell,\rho}^{-1}\xi^2}v_{\tau_0}(X)]^2
\Biggl(
\left\lvert
\sfS_\rho(\sfL(1/\rho)^{\iota_1}\xi^2)
-\sfS_\rho(\kappa_{\ell,\rho}\xi^2)
\right\rvert^{1/2}
\nonumber\\
&\hspace{16em}
+\frac{\sfL(2)^{1/2}K_{\lambda,\rho,\xi^2}^{1/2}}
{\sfL(1/\rho)^{1/2}}
\Biggr)
\nonumber\\
&\hspace{5em}
+\left(
\EE\left\lvert
\clG_{\kappa_{\ell,\rho}^{-1}\xi^2}v_{\tau_0}(X)
-\clG_{\sfL(1/\rho)^{\iota_2}\xi^2}v_t(X)
\right\rvert^2
\right)^{1/2}
\Biggr].
\tag{5.58}\label{eq:applytimeregularity}
\end{align}
Recalling the definition of $\sfS_\rho$, for $\rho$ sufficiently small,
\begin{multline}
0\leq
\sfS_\rho(\sfL(1/\rho)^{\iota_1}\xi^2)
-\sfS_\rho(\kappa_{\ell,\rho}\xi^2)
\leq
\sfS_\rho(\sfL(1/\rho)\xi^2)
-\sfS_\rho(\kappa_{\ell,\rho}\xi^2)
\\
=\frac{1}{\sfL(1/\rho)}
\log\frac{1+\sfL(1/\rho)\rho^{-1}\xi^2}
{1+\kappa_{\ell,\rho}\rho^{-1}\xi^2}
\leq
\frac{\log[\sfL(1/\rho)\kappa_{\ell,\rho}^{-1}]}
{\sfL(1/\rho)}
\leq
\frac{C\sfL(\sfL(1/\rho))}{\sfL(1/\rho)}.
\tag{5.59}\label{eq:Srhodiff}
\end{multline}
Here $C$ depends only on $\ell$. The penultimate inequality holds because
$\kappa_{\ell,\rho}<\sfL(1/\rho)$ when $\rho$ is sufficiently small.

Also, using the continuity estimate and
$\kappa_{\ell,\rho}^{-1}>\sfL(1/\rho)$, there is an absolute constant
$C<\infty$ such that, for sufficiently small $\rho$,
\begin{align}
&\left(
\EE\left\lvert
\clG_{\kappa_{\ell,\rho}^{-1}\xi^2}v_{\tau_0}(X)
-\clG_{\sfL(1/\rho)^{\iota_2}\xi^2}v_t(X)
\right\rvert^2
\right)^{1/2}
\nonumber\\
&\quad\leq C\Xi_{2,t}[v_0]
\Biggl[
\sfL_\rho\!\left(
\frac{
\kappa_{\ell,\rho}^{-1}-\sfL(1/\rho)^{\iota_2}
+\kappa_{\ell,\rho}^{-1}+\kappa_{\ell,\rho}}
{\sfL(1/\rho)^{\iota_2}}
\right)^{1/2}
\nonumber\\
&\hspace{8em}
+\left(
\frac{
[\kappa_{\ell,\rho}^{-1}-\sfL(1/\rho)^{\iota_2}
+\kappa_{\ell,\rho}^{-1}+\kappa_{\ell,\rho}]\xi^2}
{\tau_0}
\right)^{1/2}
+\left(
\frac{\sfL(\sfL(1/\rho))}{\sfL(1/\rho)}
\right)^{1/2}
\Biggr]
\nonumber\\
&\quad\overset{\eqref{eq:tauestimates}}{\leq}
C\Xi_{2,t}[v_0]
\left[
\sfL_\rho(2\kappa_{\ell,\rho}^{-1})^{1/2}
+\frac{1}{\sfL(1/\rho)^{1/2}}
+\left(
\frac{\sfL(\sfL(1/\rho))}{\sfL(1/\rho)}
\right)^{1/2}
\right]
\nonumber\\
&\quad\leq C\Xi_{2,t}[v_0]
\left[
\frac{\sfL(\sfL(1/\rho))}{\sfL(1/\rho)}
\right]^{1/2}.
\tag{5.60}\label{eq:EGdiff}
\end{align}
Here $C$ has been allowed to change from expression to expression.

Using \eqref{eq:Srhodiff} and \eqref{eq:EGdiff} in
\eqref{eq:applytimeregularity}, we obtain
\begin{align}
&\left(\EE\left\lvert
J_{\sigma,\rho}\!\left(
\sfS_\rho(\kappa_{\ell,\rho}\xi^2),
\clG_{\kappa_{\ell,\rho}^{-1}\xi^2}v_{\tau_0}(X)
\right)
-J_{\sigma,\rho}\!\left(
\sfS_\rho(\sfL(1/\rho)^{\iota_1}\xi^2),
\clG_{\sfL(1/\rho)^{\iota_2}\xi^2}v_{\tau_0}(X)
\right)
\right\rvert_{\Frob}^2\right)^{1/2}
\nonumber\\
&\quad\leq C\left(
\Xi_{2;\sfL(1/\rho)\xi^2}
[\clG_{\tau_0-\tau_1}v_{\tau_0}(X)]
+\Xi_{2;t}[v_0]
\right)
\left[
\frac{\sfL(\sfL(1/\rho))}{\sfL(1/\rho)}
\right]^{1/2}.
\tag{5.61}\label{eq:Jcontbd}
\end{align}
The constant $C$ depends on $\ell$, $\lambda$, and $\overline q$.

\medskip
\noindent\textit{Step 5: Conclusion.}
Recall that
$0<\tau_0<\tau_1<t\leq\sfT_\rho(\overline\beta^{-2})$.
Thus, by the moment bound and the hypothesis
$\overline\beta>\eta(\beta,\ell)$,
\begin{align}
&\Xi_{2;\tau_0,\tau_1}[v_{\tau_0}]^2
+\Xi_{\ell;\tau_1,\tau_2}
[\clG_{\tau_1-\tau_0}v_{\tau_0}(X)]^2
+\Xi_{2;\sfL(1/\rho)\xi^2}
[\clG_{\tau_0-\tau_1}v_{\tau_0}(X)]
+\Xi_{2;t}[v_0]
\nonumber\\
&\qquad\leq C\left\langle\triplenorm{v_0}_\ell\right\rangle^2,
\tag{5.62}\label{eq:applymomentbd}
\end{align}
where the constant $C$ depends on $M,\beta,\overline\beta,\ell$.

Combining \eqref{eq:E1OK}, \eqref{eq:E2finalbd}, and
\eqref{eq:Jcontbd} using the triangle inequality, Cauchy--Schwarz, and
\eqref{eq:applymomentbd}, we see that
\begin{align*}
\EE\left\lvert
\left(\fint_{\ttR}\sigma^2(v_t(y))\,\dif y\right)^{1/2}
-J_{\sigma,\rho}\!\left(
\sfS_\rho(\sfL(1/\rho)^{\iota_1}\xi^2),
\clG_{\sfL(1/\rho)^{\iota_2}\xi^2}v_{\tau_0}(x)
\right)
\right\rvert_{\Frob}^{2}
\leq
C\left\langle\triplenorm{v_0}_\ell\right\rangle^2
\frac{\sfL(\sfL(1/\rho))}{\sfL(1/\rho)}.
\end{align*}
The constant $C<\infty$ depends only on $q,M,\beta,\lambda,\ell$.
Recalling the definition of the square $\ttR$, this proves (5.51).
\end{proof}

\subsection{Some Volterra Estimation}

\begin{proposition}[Comparison of time-shifted solutions]
\label{prop:time-shift-comparison}

Let $0\leq \tau_0<\tau_1<\tau_2$ and $\rho\in(0,1)$.
Define
\[
L(x):=\log(1+x),
\qquad
\gamma_\rho^2:=\frac{4\pi}{L(1/\rho)}.
\]

Let
\[
G_t(x)
=
\frac{1}{2\pi t}
\exp\left(-\frac{|x|^2}{2t}\right)
\]
be the heat kernel on $\mathbb{R}^2$.

Suppose that $\sigma_\rho:\mathbb{R}\to\mathbb{R}$
satisfies
\[
\|\sigma_\rho\|_\infty\leq M,
\qquad
\operatorname{Lip}(\sigma_\rho)
\leq \gamma_\rho^{-1}\lambda,
\]
where $M$ is independent of $\rho$.

Let $v^\rho$ and $\widetilde v^\rho$ be
two adapted mild solutions driven by the same
space-time white noise $W$, satisfying
\begin{align*}
v_t^\rho(y)
={}&
\bigl(G_{t-\tau_0}*v_{\tau_0}^\rho\bigr)(y)
\\
&+
\gamma_\rho
\int_{\tau_0}^{t}
\int_{\mathbb{R}^2}
G_{t+\rho-r}(y-x)
\sigma_\rho(v_r^\rho(x))
\,W(dr,dx),
\end{align*}
and
\begin{align*}
\widetilde v_t^\rho(y)
={}&
\bigl(
G_{t-\tau_1}*
G_{\tau_1-\tau_0}*
v_{\tau_0}^\rho
\bigr)(y)
\\
&+
\gamma_\rho
\int_{\tau_1}^{t}
\int_{\mathbb{R}^2}
G_{t+\rho-r}(y-x)
\sigma_\rho(\widetilde v_r^\rho(x))
\,W(dr,dx).
\end{align*}

Define
\[
F(t)
:=
\sup_{y\in\mathbb{R}^2}
\mathbb{E}
\left[
|v_t^\rho(y)-\widetilde v_t^\rho(y)|^2
\right].
\]

Then, for every $t\in[\tau_1,\tau_2]$,
\begin{equation}
\label{eq:time-shift-volterra}
{
\begin{aligned}
F(t)
\leq{}&
\frac{M^2}{L(1/\rho)}
L\left(
\frac{\tau_1-\tau_0}
{t-\tau_1+\rho}
\right) +
\frac{1}{4\pi}
\int_{\tau_1}^{t}
\frac{F(r)}{t+\rho-r}\,dr.
\end{aligned}
}
\end{equation}

Moreover, the following explicit bound holds:
\begin{equation}
\label{eq:time-shift-explicit}
{
\begin{aligned}
F(t)
\leq
\frac{M^2}{L(1/\rho)}
\bigg[
&
L\left(
\frac{\tau_1-\tau_0}
{t-\tau_1+\rho}
\right) + 
4L\left(
\frac{t-\tau_1}{\rho}
\right)
\bigg].
\end{aligned}
}
\end{equation}

\end{proposition}


\begin{proof}

By the semigroup property of the heat kernel,
\[
G_{t-\tau_1}*G_{\tau_1-\tau_0}
=
G_{t-\tau_0}.
\]

Consequently, for $t\geq\tau_1$,
\begin{align*}
&v_t^\rho(y)-\widetilde v_t^\rho(y)
\\
={}&
\gamma_\rho
\int_{\tau_0}^{\tau_1}
\int_{\mathbb{R}^2}
G_{t+\rho-r}(y-x)
\sigma_\rho(v_r^\rho(x))
\,W(dr,dx)
\\
&+
\gamma_\rho
\int_{\tau_1}^{t}
\int_{\mathbb{R}^2}
G_{t+\rho-r}(y-x)
\\
&\qquad\qquad\times
\left[
\sigma_\rho(v_r^\rho(x))
-
\sigma_\rho(\widetilde v_r^\rho(x))
\right]
\,W(dr,dx).
\end{align*}

Since the two stochastic integrals are taken
over disjoint time intervals, they are orthogonal
in $L^2(\Omega)$.

Applying the It\^o isometry gives
\begin{align*}
&\mathbb{E}
\left[
|v_t^\rho(y)-\widetilde v_t^\rho(y)|^2
\right]
\nonumber\\
={}&
\gamma_\rho^2
\int_{\tau_0}^{\tau_1}
\int_{\mathbb{R}^2}
G_{t+\rho-r}^2(y-x)
\mathbb{E}
\left[
|\sigma_\rho(v_r^\rho(x))|^2
\right]
\,dx\,dr
\nonumber\\
&+
\gamma_\rho^2
\int_{\tau_1}^{t}
\int_{\mathbb{R}^2}
G_{t+\rho-r}^2(y-x)
\mathbb{E}
\left[
\left|
\sigma_\rho(v_r^\rho(x))
-
\sigma_\rho(\widetilde v_r^\rho(x))
\right|^2
\right]
\,dx\,dr.
\end{align*}

Recall that in two dimensions,
\[
G_t^2(x)
=
\frac{1}{4\pi t}G_{t/2}(x).
\]

In particular,
\[
\int_{\mathbb{R}^2}G_t^2(x)\,dx
=
\frac{1}{4\pi t}.
\]

Using the boundedness of $\sigma_\rho$,
we have
\begin{align*}
&\gamma_\rho^2 M^2
\int_{\tau_0}^{\tau_1}
\int_{\mathbb{R}^2}
G_{t+\rho-r}^2(y-x)
\,dx\,dr
\\
={}&
\frac{\gamma_\rho^2M^2}{4\pi}
\int_{\tau_0}^{\tau_1}
\frac{1}{t+\rho-r}\,dr
\\
={}&
\frac{M^2}{L(1/\rho)}
\log\left(
\frac{t-\tau_0+\rho}
{t-\tau_1+\rho}
\right)
\\
={}&
\frac{M^2}{L(1/\rho)}
L\left(
\frac{\tau_1-\tau_0}
{t-\tau_1+\rho}
\right).
\end{align*}

For the second term, the Lipschitz assumption
implies
\[
\begin{aligned}
&\left|
\sigma_\rho(v_r^\rho(x))
-
\sigma_\rho(\widetilde v_r^\rho(x))
\right|^2 \leq
\gamma_\rho^{-2}\lambda^2
|v_r^\rho(x)-\widetilde v_r^\rho(x)|^2.
\end{aligned}
\]

Therefore,
\begin{align*}
&\gamma_\rho^2
\int_{\tau_1}^{t}
\int_{\mathbb{R}^2}
G_{t+\rho-r}^2(y-x)
\mathbb{E}
\left[
\left|
\sigma_\rho(v_r^\rho(x))
-
\sigma_\rho(\widetilde v_r^\rho(x))
\right|^2
\right]
\,dx\,dr
\\
\leq{}&\lambda^2
\int_{\tau_1}^{t}
\int_{\mathbb{R}^2}
G_{t+\rho-r}^2(y-x)
F(r)
\,dx\,dr
\\
={}&
\frac{\lambda^2}{4\pi}
\int_{\tau_1}^{t}
\frac{F(r)}{t+\rho-r}\,dr.
\end{align*}

Combining the two estimates and taking
the supremum over $y\in\mathbb{R}^2$
establishes \eqref{eq:time-shift-volterra}.

Finally, since
\[
|\sigma_\rho(a)-\sigma_\rho(b)|
\leq 2M,
\]
we may directly estimate the second term
in \eqref{eq:time-shift-ito} by
\begin{align*}
&4\gamma_\rho^2M^2
\int_{\tau_1}^{t}
\int_{\mathbb{R}^2}
G_{t+\rho-r}^2(y-x)
\,dx\,dr
\\
={}&
\frac{4M^2}{L(1/\rho)}
\int_{\tau_1}^{t}
\frac{1}{t+\rho-r}\,dr
\\
={}&
\frac{4M^2}{L(1/\rho)}
L\left(
\frac{t-\tau_1}{\rho}
\right).
\end{align*}

Together with the estimate for the first
term, this proves
\eqref{eq:time-shift-explicit}.

\end{proof}


\begin{remark}[The associated Volterra inequality]
\label{rem:time-shift-volterra}

Set
\[
\Delta:=\tau_1-\tau_0,
\qquad
g(\tau):=F(\tau_1+\tau).
\]

Then \eqref{eq:time-shift-volterra}
takes the form
\begin{equation}
\label{eq:shifted-volterra}
g(\tau)
\leq
A_\rho
L\left(
\frac{\Delta}{\tau+\rho}
\right)
+
B\int_0^\tau
\frac{g(r)}{\tau+\rho-r}\,dr,
\end{equation}
where
\[
A_\rho=\frac{M^2}{L(1/\rho)},
\qquad
B=\frac{1}{4\pi}.
\]

To estimate the convolution term, define
\[
f_\rho(\tau)
:=
L\left(
\frac{\Delta}{\tau+\rho}
\right).
\]

For $0\leq r\leq\tau$, we have
\begin{align*}
f_\rho(r)-f_\rho(\tau)
={}&
\log\left(
\frac{r+\rho+\Delta}
{\tau+\rho+\Delta}
\right)
\\
&+
\log\left(
\frac{\tau+\rho}{r+\rho}
\right)
\\
\leq{}&
\log\left(
\frac{\tau+\rho}{r+\rho}
\right).
\end{align*}

Consequently,
\begin{align*}
&\int_0^\tau
\frac{f_\rho(r)}{\tau+\rho-r}\,dr
\\
\leq{}&
f_\rho(\tau)
\int_0^\tau
\frac{1}{\tau+\rho-r}\,dr
\\
&+
\int_0^\tau
\frac{
\log((\tau+\rho)/(r+\rho))
}{\tau+\rho-r}\,dr
\\
\leq{}&
f_\rho(\tau)L(\tau/\rho)
+
\int_0^\tau
\frac{\log(\tau/r)}{\tau-r}\,dr
\\
={}&
f_\rho(\tau)L(\tau/\rho)
+
\int_0^1
\frac{-\log s}{1-s}\,ds
\\
={}&
f_\rho(\tau)L(\tau/\rho)
+\frac{\pi^2}{6}.
\end{align*}

Thus,
\begin{equation}
\label{eq:log-convolution-estimate}
\boxed{
\begin{aligned}
&\int_0^\tau
\frac{
L\left(\frac{\Delta}{r+\rho}\right)
}{\tau+\rho-r}\,dr
\\
&\qquad\leq
L\left(
\frac{\Delta}{\tau+\rho}
\right)
L\left(
\frac{\tau}{\rho}
\right)
+
\frac{\pi^2}{6}.
\end{aligned}
}
\end{equation}

However, the coefficient $B$ in
\eqref{eq:shifted-volterra} is independent
of $\rho$, and
\[
B L(T/\rho)\longrightarrow\infty
\qquad\text{as }\rho\downarrow0
\]
for every fixed $T>0$.

Hence a direct contraction argument for
\eqref{eq:shifted-volterra} on a fixed
time interval does not provide a
$\rho$-uniform refined bound.

\end{remark}
